\documentclass[11pt,a4paper]{article}
\usepackage{amsmath,amssymb,amsthm}
\usepackage{mathtools}
\usepackage{bm}
\usepackage{booktabs}
\usepackage{array}
\usepackage{enumitem}
\usepackage{geometry}
\usepackage{xcolor}
\usepackage{tcolorbox}
\usepackage{listings}
\usepackage[pdfencoding=auto,psdextra]{hyperref}

\geometry{margin=2.5cm}

\newcommand{\E}{\mathbb{E}}
\newcommand{\R}{\mathbb{R}}
\newcommand{\norm}[1]{\left\lVert #1 \right\rVert}
\newcommand{\dt}{\Delta t}
\newcommand{\dk}{d_k}

\newtheorem*{remark*}{Remark}

\lstset{
  basicstyle=\ttfamily\small,
  breaklines=true,
  columns=fullflexible,
  keepspaces=true,
  frame=single,
  framesep=4pt,
  xleftmargin=4pt,
}

\title{\textbf{Why late timesteps dominate the XOPD per-step loss}\\[0.3em]
\large A flow-matching account of the \texttt{d\_k} blow-up, artifact vs.\ signal,
and the \texttt{xopd\_dk\_space} reweighting}
\author{Flow Factory}
\date{}

\begin{document}
\maketitle

\begin{abstract}
Theory companion to the empirical observation
\texttt{docs/xopd/progress/2026-07-12-per-timestep-dk-dominance.md} (last quarter of steps
$\approx 81\%$ of $\sum_k \dk$, final step $\approx 40\%$); the implemented remedy is
\texttt{train.xopd\_dk\_space="x0"} (clean-latent loss, \S\ref{sec:impl}). Relevant
code: \texttt{compute\_per\_step\_kl} in
\texttt{src/flow\_factory/trainers/xopd/common.py}, called by
\texttt{\_precompute\_d\_per\_timestep} / \texttt{\_optimize\_train\_pass} in
\texttt{src/flow\_factory/trainers/xopd/trainer.py}.

This note defines three per-step MSE losses ($L_{xt}$, $L_v$, $L_{x_0}$), their exact conversion
on the FLUX.2-klein 28-step grid, and why the default $L_{xt}$ (config \texttt{"xt"}) dominates late
timesteps. Companion data: \S\ref{sec:schedule}--\S\ref{sec:empirical}.
\end{abstract}

\tableofcontents

%% ============================================================
\section{Three per-step MSE losses and their exact conversion}
\label{sec:three-losses}
%% ============================================================

XOPD compares student and teacher on the \emph{same} on-policy rollout state $x_t$ at step
index~$k$; only their predictions differ. Under rectified-flow interpolation
($x_t=(1-\sigma)z_0+\sigma\varepsilon$, $\sigma=t/1000$) and ODE Euler stepping
($\mu=x_t+v\,\dt$, $\dt=\sigma^{\mathrm{sched}}_{k+1}-\sigma^{\mathrm{sched}}_k<0$), there are
three equivalent MSE targets (config \texttt{xopd\_dk\_space}, default \texttt{"xt"}):

\begin{center}
\renewcommand{\arraystretch}{1.25}
\begin{tabular}{@{}clll@{}}
\toprule
\textbf{knob} & \textbf{loss} & \textbf{formula} & \textbf{note} \\
\midrule
\texttt{xt} & $L_{xt}$ & $\norm{\mu_s-\mu_t}^2$ & transition mean / next latent (default) \\
\texttt{v} & $L_v$ & $\norm{v_s-v_t}^2$ & raw velocity MSE; ODE-only \\
\texttt{x0} & $L_{x_0}$ & $\norm{z_{0,s}-z_{0,t}}^2$ & clean-latent MSE; ODE-only \\
\bottomrule
\end{tabular}
\end{center}

\textbf{Ratio.} $\mathrm{MSE}(v):\mathrm{MSE}(xt):\mathrm{MSE}(x0) = 1:\dt^2:\sigma^2$.

Here $\mu$ is the one-step transition mean ($x_{t+\dt}$ under ODE), $v$ the model velocity
(\texttt{noise\_pred}), and $z_0=x_t-\sigma v$ the analytically recovered clean latent. Because
$x_t$ cancels:
\begin{align}
\mu_s-\mu_t &= \dt\,(v_s-v_t),
\label{eq:mu-v}\\
z_{0,s}-z_{0,t} &= -\sigma\,(v_s-v_t),
\label{eq:x0-v}\\
\Rightarrow\quad
L_{xt} &= \dt^2\,L_v,
\label{eq:mu-v-sq}\\
L_{x_0} &= \sigma^2\,L_v,
\label{eq:x0-v-sq}\\
L_{x_0} &= \Big(\tfrac{\sigma}{\dt}\Big)^2 L_{xt},
\label{eq:x0-mu}\\
L_v &= \Big(\tfrac{\dt}{\sigma}\Big)^2 L_{x_0}.
\label{eq:v-x0}
\end{align}
All four identities are \emph{exact} on the discrete ODE grid (no approximation). The conversion
factors depend \emph{only} on $(\sigma_k,\dt_k)$ at step~$k$; they do not require equal errors
across steps --- they convert \emph{the same} underlying $(v_s-v_t)$ disagreement between spaces.

\paragraph{Conversion matrix (multiply to go from column loss to row loss).}
\begin{center}
\footnotesize
\renewcommand{\arraystretch}{1.3}
\begin{tabular}{@{}l|ccc@{}}
 & $\to L_{xt}$ & $\to L_v$ & $\to L_{x_0}$ \\
\midrule
from $L_{xt}$ & $1$ & $\dt^{-2}$ & $(\sigma/\dt)^2$ \\
from $L_v$ & $\dt^2$ & $1$ & $\sigma^2$ \\
from $L_{x_0}$ & $(\dt/\sigma)^2$ & $\sigma^{-2}$ & $1$ \\
\end{tabular}
\end{center}

\paragraph{What each loss implicitly weights (if you sum $\sum_k L_{\{\cdot\},k}$ with coefficient~$1$).}
Fix the \emph{underlying} velocity disagreement $L_{v,k}$ at step~$k$:
\begin{itemize}[leftmargin=1.5em]
\item \textbf{$L_{xt}$ (default):} step~$k$ contributes $\dt_k^2 L_{v,k}$ --- late steps have larger
$\lvert\dt_k\rvert$ and dominate (\S\ref{sec:schedule}, col.~``\%$L_{xt}$'').
\item \textbf{$L_v$ (\texttt{xopd\_dk\_space="v"}):} step~$k$ contributes $L_{v,k}$ directly --- the
native FM velocity target; under the fixed 28-step grid the per-step $\dt_k$ vary only modestly so
this is close to uniform (\S\ref{sec:empirical}).
\item \textbf{$L_{x_0}$ (\texttt{xopd\_dk\_space="x0"}):} step~$k$ contributes $\sigma_k^2 L_{v,k}$ ---
early high-noise steps ($\sigma\approx1$) dominate (col.~``pct $L_{x_0}$'').
\end{itemize}
Alternatively, fix $L_{x_0,k}$ (equal clean-latent disagreement): $L_{xt,k}=(\dt_k/\sigma_k)^2
L_{x_0,k}$, which is the ``$(\dt/\sigma)^2$ blow-up'' narrative (\S\ref{sec:causes}).

%% ============================================================
\section{The key identity (legacy shorthand)}
\label{sec:identity}
%% ============================================================

With \texttt{normalize\_d\_k=false}, the logged \texttt{train/d\_k/\{ti\}} is $L_{xt}$ (not $L_v$).
Combining \eqref{eq:mu-v-sq}--\eqref{eq:v-x0}:
\[
L_{xt} = \dt^2 L_v = \Big(\tfrac{\dt}{\sigma}\Big)^2 L_{x_0}.
\]
So the default loss is clean-latent disagreement multiplied by the implicit per-step weight
$(\dt/\sigma)^2$ (\S\ref{sec:three-losses}). The $x$-space knob sets $L_{x_0}$ directly and removes
that factor (\S\ref{sec:impl}).

%% ============================================================
\section{Concrete schedule: official FLUX.2-klein @ 28 steps}
\label{sec:schedule}
%% ============================================================

We reproduce the \emph{exact} denoising grid used in XOPD training
(\texttt{num\_inference\_steps=28}, 512$\times$512) by loading the default scheduler from the
official HuggingFace repo and calling the same \texttt{retrieve\_timesteps} path as
\texttt{Flux2KleinPipeline} (shifted flow-matching sigmas with
\texttt{image\_seq\_len=1024}, $\mu=1.859$):
\begin{lstlisting}[language=Python,basicstyle=\ttfamily\footnotesize,frame=none]
FlowMatchEulerDiscreteScheduler.from_pretrained(
    "black-forest-labs/FLUX.2-klein-base-4B", subfolder="scheduler")
# sigmas = linspace(1, 1/28, 28); mu = compute_empirical_mu(1024, 28)
\end{lstlisting}
Repro script: \texttt{.scratch/flux2\_klein\_scheduler\_timesteps.py} (uv venv
\texttt{/Users/bowenping/.venvs/flux2-klein-scheduler}, \texttt{diffusers==0.39.0}).

\paragraph{Notation at step index \texttt{ti}.}
\texttt{ti=0} is the noisiest state; \texttt{ti=27} is the last denoising step before $\sigma=0$.
The scheduler returns a \emph{timestep} $t_k\in[0,1000]$ and a \emph{shifted flow sigma}
$\sigma^{\mathrm{sched}}_k\in[0,1]$. XOPD's clean-latent recovery uses the repo's FM fraction
$\sigma_k=t_k/1000$ (\texttt{\_noise\_fraction}, \texttt{flow\_match\_sigma}). The ODE step size is
$\Delta t_k=\sigma^{\mathrm{sched}}_{k+1}-\sigma^{\mathrm{sched}}_k<0$ (from
\texttt{scheduler.step}). On this grid $\sigma_k\approx\sigma^{\mathrm{sched}}_k$ to machine
precision.

\paragraph{Per-step conversion factors on the 28-step grid.}
Table~\ref{tab:klein28-summary} lists representative steps; Table~\ref{tab:klein28-full}
(auto-generated, \texttt{docs/xopd/flux2\_klein\_28step\_loss\_scales.tex}) gives all 28 rows.
Repro: \texttt{.scratch/flux2\_klein\_scheduler\_timesteps.py}.

\begin{table}[ht]
\centering
\footnotesize
\renewcommand{\arraystretch}{1.15}
\begin{tabular}{@{}rrrrrrrrr@{}}
\toprule
\texttt{ti} & $t$ & $\sigma$ & $\Delta t$
 & $\frac{L_{xt}}{L_v}=\dt^2$ & $\frac{L_{x_0}}{L_v}=\sigma^2$ & $\frac{L_{x_0}}{L_{xt}}$
 & pct $L_{xt}$ & pct $L_{x_0}$ \\
\midrule
 0 & 1000.0 & 1.0000 & $-$0.00574 & $3.3{\times}10^{-5}$ & 1.000 & $3.0{\times}10^{4}$ & 0.04\% & 5.28\% \\
 7 &  950.6 & 0.9506 & $-$0.00929 & $8.6{\times}10^{-5}$ & 0.904 & $1.0{\times}10^{4}$ & 0.10\% & 4.77\% \\
13 &  881.0 & 0.8810 & $-$0.01584 & $2.5{\times}10^{-4}$ & 0.776 & $3.1{\times}10^{3}$ & 0.28\% & 4.10\% \\
21 &  681.5 & 0.6815 & $-$0.04505 & $2.0{\times}10^{-3}$ & 0.464 & $2.3{\times}10^{2}$ & 2.27\% & 2.45\% \\
24 &  516.8 & 0.5168 & $-$0.08175 & $6.7{\times}10^{-3}$ & 0.267 & 40.0 & 7.47\% & 1.41\% \\
25 &  435.1 & 0.4351 & $-$0.10456 & $1.1{\times}10^{-2}$ & 0.189 & 17.3 & 12.21\% & 1.00\% \\
26 &  330.5 & 0.3305 & $-$0.13847 & $1.9{\times}10^{-2}$ & 0.109 & 5.70 & 21.42\% & 0.58\% \\
27 &  192.1 & 0.1921 & $-$0.19206 & $3.7{\times}10^{-2}$ & 0.037 & 1.00 & 41.21\% & 0.19\% \\
\bottomrule
\end{tabular}
\caption{FLUX.2-klein-base-4B, \texttt{num\_inference\_steps=28}, 512px. Columns 5--7:
exact factors from \eqref{eq:mu-v-sq}--\eqref{eq:x0-mu}. ``\%$L_{xt}$'' (``\%$L_{x_0}$''): share
of $\sum_k L_{xt,k}$ (resp.\ $\sum_k L_{x_0,k}$) if $L_{v,k}\equiv$ const across steps.}
\label{tab:klein28-summary}
\end{table}

% Auto-generated by .scratch/flux2_klein_scheduler_timesteps.py — do not edit by hand.
\begin{table}[ht]
\centering
\scriptsize
\setlength{\tabcolsep}{3pt}
\renewcommand{\arraystretch}{1.05}
\begin{tabular}{@{}rrrrrrrrr@{}}
\toprule
\texttt{ti} & $t$ & $\sigma$ & $\Delta t$
 & $\frac{L_{xt}}{L_v}$ & $\frac{L_{x_0}}{L_v}$ & $\frac{L_{x_0}}{L_{xt}}$
 & pct $L_{xt}$ & pct $L_{x_0}$ \\
\midrule
 0 &  1000.0 & 1.0000 & -0.00574 & 3.3e-05 & 1.000 & 30380.38 & 0.04\% & 5.28\% \\
 1 &   994.3 & 0.9943 & -0.00611 & 3.7e-05 & 9.89e-01 & 26518.80 & 0.04\% & 5.22\% \\
 2 &   988.2 & 0.9882 & -0.00651 & 4.2e-05 & 9.76e-01 & 23039.00 & 0.05\% & 5.15\% \\
 3 &   981.6 & 0.9816 & -0.00696 & 4.8e-05 & 9.64e-01 & 19912.21 & 0.05\% & 5.09\% \\
 4 &   974.7 & 0.9747 & -0.00745 & 5.6e-05 & 9.50e-01 & 17114.06 & 0.06\% & 5.01\% \\
 5 &   967.2 & 0.9672 & -0.00800 & 6.4e-05 & 9.36e-01 & 14621.44 & 0.07\% & 4.94\% \\
 6 &   959.2 & 0.9592 & -0.00861 & 7.4e-05 & 9.20e-01 & 12410.80 & 0.08\% & 4.86\% \\
 7 &   950.6 & 0.9506 & -0.00929 & 8.6e-05 & 9.04e-01 & 10460.53 & 0.10\% & 4.77\% \\
 8 &   941.3 & 0.9413 & -0.01006 & 1.0e-04 & 8.86e-01 & 8748.98 & 0.11\% & 4.68\% \\
 9 &   931.3 & 0.9313 & -0.01093 & 1.2e-04 & 8.67e-01 & 7256.11 & 0.13\% & 4.58\% \\
10 &   920.3 & 0.9203 & -0.01192 & 1.4e-04 & 8.47e-01 & 5962.25 & 0.16\% & 4.47\% \\
11 &   908.4 & 0.9084 & -0.01305 & 1.7e-04 & 8.25e-01 & 4849.23 & 0.19\% & 4.35\% \\
12 &   895.4 & 0.8954 & -0.01434 & 2.1e-04 & 8.02e-01 & 3899.25 & 0.23\% & 4.23\% \\
13 &   881.0 & 0.8810 & -0.01584 & 2.5e-04 & 7.76e-01 & 3095.63 & 0.28\% & 4.10\% \\
14 &   865.2 & 0.8652 & -0.01758 & 3.1e-04 & 7.49e-01 & 2422.59 & 0.35\% & 3.95\% \\
15 &   847.6 & 0.8476 & -0.01963 & 3.9e-04 & 7.18e-01 & 1865.26 & 0.43\% & 3.79\% \\
16 &   828.0 & 0.8280 & -0.02205 & 4.9e-04 & 6.86e-01 & 1409.55 & 0.54\% & 3.62\% \\
17 &   805.9 & 0.8059 & -0.02496 & 6.2e-04 & 6.50e-01 & 1042.44 & 0.70\% & 3.43\% \\
18 &   781.0 & 0.7810 & -0.02849 & 8.1e-04 & 6.10e-01 & 751.66 & 0.91\% & 3.22\% \\
19 &   752.5 & 0.7525 & -0.03281 & 1.08e-03 & 5.66e-01 & 525.93 & 1.20\% & 2.99\% \\
20 &   719.7 & 0.7197 & -0.03821 & 1.46e-03 & 5.18e-01 & 354.83 & 1.63\% & 2.73\% \\
21 &   681.5 & 0.6815 & -0.04505 & 2.03e-03 & 4.64e-01 & 228.85 & 2.27\% & 2.45\% \\
22 &   636.4 & 0.6364 & -0.05391 & 2.91e-03 & 4.05e-01 & 139.37 & 3.25\% & 2.14\% \\
23 &   582.5 & 0.5825 & -0.06567 & 4.31e-03 & 3.39e-01 & 78.683 & 4.82\% & 1.79\% \\
24 &   516.8 & 0.5168 & -0.08175 & 6.68e-03 & 2.67e-01 & 39.970 & 7.47\% & 1.41\% \\
25 &   435.1 & 0.4351 & -0.10456 & 1.09e-02 & 1.89e-01 & 17.315 & 12.21\% & 1.00\% \\
26 &   330.5 & 0.3305 & -0.13847 & 1.92e-02 & 1.09e-01 & 5.698 & 21.42\% & 0.58\% \\
27 &   192.1 & 0.1921 & -0.19206 & 3.69e-02 & 3.69e-02 & 1.000 & 41.21\% & 0.19\% \\
\bottomrule
\end{tabular}
\caption{Three per-step MSE losses on the FLUX.2-klein 28-step grid ($512\times512$).
Columns 5--7 are exact conversion factors at step $k$ (on-policy, shared $x_t$).
Last two columns: normalized share of $\sum_k L_{xt}$ (resp.\ $\sum_k L_{x_0}$)
if $\norm{v_s-v_t}^2$ were identical at every step.
Generated from \texttt{black-forest-labs/FLUX.2-klein-base-4B} default scheduler.}
\label{tab:klein28-full}
\end{table}


\paragraph{Reading the numbers.}
\begin{itemize}[leftmargin=1.5em]
\item \textbf{At \texttt{ti=27}, all three losses coincide in scale.} $\lvert\dt_{27}\rvert\approx
\sigma_{27}$, so $\dt_{27}^2=\sigma_{27}^2\approx0.037$ and $L_{xt}=L_{x_0}=L_v$ for the \emph{same}
velocity gap.
\item \textbf{Default $L_{xt}$ dominance $\approx$ flat $L_v$.} If $\norm{v_s-v_t}^2$ were identical
at every step, $\sum_k L_{xt,k}$ puts $\mathbf{41\%}$ on \texttt{ti=27} (col.~``\%$L_{xt}$'').
The mixed XOPD run observes $\mathbf{40\%}$ on \texttt{ti=27} --- the logged curve matches
\emph{roughly constant velocity error} amplified only by $\dt_k^2$, not the stronger
$(\dt_k/\sigma_k)^2$ factor (which would give $\mathbf{77\%}$ if $L_{x_0,k}$ were constant).
\item \textbf{$L_{x_0}$ up-weights early steps.} With flat $L_v$, high-noise steps each take
$\sim5\%$ of $\sum L_{x_0}$ while \texttt{ti=27} takes only $\mathbf{0.19\%}$. Setting
\texttt{xopd\_dk\_space="x0"} switches to this allocation.
\item \textbf{Cross-step ratios (same $L_v$).} $L_{xt,27}/L_{xt,13}=(\dt_{27}/\dt_{13})^2
\approx\mathbf{148\times}$; $L_{x_0,0}/L_{x_0,27}=\sigma_0^2/\sigma_{27}^2\approx\mathbf{27\times}$.
\item \textbf{``Late step'' $\neq$ $\sigma\to0$.} The last rollout step has $t=192$ ($\sigma=0.19$),
not $t\approx0$. Dominance comes from the discrete grid's large $\lvert\dt_k\rvert$ at the end.
\end{itemize}

%% ============================================================
\section{Empirical distribution: mixed run, converted from measured $L_{xt}$}
\label{sec:empirical}
%% ============================================================

Wandb run \texttt{4cnkluwk} (mixed geneval+OCR, \texttt{xopd\_dk\_space=xt},
\texttt{normalize\_d\_k=false}) logs $L_{xt,k}$ directly
(\S\ref{sec:causes}). Under ODE we convert each step to the \emph{hypothetical} $L_v$ and $L_{x_0}$
that would have been logged with the same underlying velocity gap:
\[
L_{v,k}=\frac{L_{xt,k}}{\dt_k^2},\qquad
L_{x_0,k}=\Big(\frac{\sigma_k}{\dt_k}\Big)^2 L_{xt,k}.
\]
Table~\ref{tab:empirical-converted} uses \texttt{mean\_tail} from
\texttt{progress/2026-07-12-per-timestep-dk-dominance.md} and the official 28-step
$(\sigma_k,\dt_k)$ from \S\ref{sec:schedule}.

\begin{table}[ht]
\centering
\footnotesize
\renewcommand{\arraystretch}{1.12}
\begin{tabular}{@{}lrrr@{}}
\toprule
\textbf{aggregate} & \textbf{$L_{xt}$ (meas.)} & \textbf{$L_v$ (conv.)} & \textbf{$L_{x_0}$ (conv.)} \\
\midrule
top-1 (\texttt{ti=27}) & 40.2\% & 0.5\% & 0.02\% \\
top-7 (\texttt{ti=21--27}, last 1/4) & 81.0\% & 3.1\% & 0.9\% \\
high-noise half (\texttt{ti=0--13}) & 12.2\% & 91.5\% & 95.2\% \\
low-noise half (\texttt{ti=14--27}) & 87.8\% & 8.5\% & 4.8\% \\
\bottomrule
\end{tabular}
\caption{Share of $\sum_k L_{\{\cdot\},k}$ (\texttt{mean\_tail}) in each loss space. $L_{xt}$ is
measured; $L_v$ and $L_{x_0}$ are ODE-exact conversions using the FLUX.2-klein 28-step grid.}
\label{tab:empirical-converted}
\end{table}

\paragraph{Reading.}
\begin{itemize}[leftmargin=1.5em]
\item \textbf{Late dominance is specific to $L_{xt}$.} The observed 40\%/81\% concentration on late
steps \emph{disappears} under conversion to $L_v$ (early half $\mathbf{92\%}$) or $L_{x_0}$
($\mathbf{95\%}$): the logged curve is dominated by $\dt_k^2$, while the underlying velocity gap
is larger at high noise.
\item \textbf{$L_v$ is nearly flat in absolute scale} after conversion ($L_{v,0}\approx1.26$ vs.\
$L_{v,27}\approx0.056$, only $\mathbf{0.5\%}$ of the sum on \texttt{ti=27}), consistent with
Cause~A (\S\ref{ssec:causeA}).
\item \textbf{Ablation implication.} \texttt{xopd\_dk\_space="v"} removes the $\dt_k^2$ rescaling
and should redistribute gradient toward high-noise steps; \texttt{"x0"} applies $\sigma_k^2$ and
amplifies that further.
\end{itemize}

%% ============================================================
\section{Implementation: the \texttt{xopd\_dk\_space} knob}
\label{sec:impl}
%% ============================================================

\textbf{Knob.} \texttt{train.xopd\_dk\_space} $\in\{\texttt{v},\texttt{xt},\texttt{x0}\}$ (default
\texttt{xt}). See \texttt{compute\_per\_step\_kl} in \texttt{xopd/common.py}; \texttt{v} and
\texttt{x0} require ODE (\texttt{XOPDTrainer.\_\_init\_\_} guard).

\begin{center}
\renewcommand{\arraystretch}{1.3}
\footnotesize
\begin{tabular}{@{}p{1.6cm}p{5.8cm}p{3.8cm}@{}}
\toprule
\textbf{knob} & \textbf{per-step loss} & \textbf{relation to $L_v$} \\
\midrule
\texttt{xt} & $L_{xt}=\norm{\mu_s-\mu_t}^2$ & $L_{xt}=\dt^2 L_v$ \\
\texttt{v} & $L_v=\norm{v_s-v_t}^2$ & identity \\
\texttt{x0} & $L_{x_0}=\norm{z_{0,s}-z_{0,t}}^2$ & $L_{x_0}=\sigma^2 L_v$ \\
\bottomrule
\end{tabular}
\end{center}
Since $x_t$ cancels, $L_{x_0}=(\sigma/\dt)^2 L_{xt}$ (\eqref{eq:x0-mu}). \texttt{normalize\_d\_k}
applies only to \texttt{xt}; \texttt{v}/\texttt{x0} ignore it.

\textbf{No extra cost.} \texttt{x0} recovers $z_0$ analytically from $\mu$; \texttt{v} uses
$(\mu_s-\mu_t)/\dt$. \textbf{Logging.} \texttt{train/d\_k/\{ti\}} values are in the selected space;
compare \emph{shapes} across spaces, not raw magnitudes (\S\ref{sec:empirical}).

%% ============================================================
\section{Two causes of the dominance: separate the artifact from the signal}
\label{sec:causes}
%% ============================================================

\subsection{Cause A --- $\dt_k^2$ amplification in $L_{xt}$ (a parameterization artifact)}
\label{ssec:causeA}

The default loss is $L_{xt}=\dt^2 L_v$ (\eqref{eq:mu-v-sq}), not velocity MSE. Summing
$\sum_k L_{xt,k}$ with unit weight therefore applies a per-step factor $\dt_k^2$ to the underlying
velocity gap. Late rollout steps have larger $\lvert\dt_k\rvert$ (Table~\ref{tab:klein28-summary}),
so \textbf{$L_{xt}$ over-represents late steps even if $L_v$ were flat}.

Quantitative check on the FLUX.2-klein 28-step grid (\S\ref{sec:schedule}):
\begin{itemize}[leftmargin=1.5em]
\item If $L_{v,k}\equiv$ const: $\sum_k L_{xt,k}$ puts $\mathbf{41\%}$ on \texttt{ti=27}
(col.~``\%$L_{xt}$''). The mixed XOPD run observes $\mathbf{40\%}$ on \texttt{ti=27} for logged
$L_{xt}$ --- the curve matches \emph{roughly constant velocity error} $\times\,\dt_k^2$, not a
separate ``mystery'' effect.
\item If instead $L_{x_0,k}\equiv$ const: $L_{xt,k}=(\dt_k/\sigma_k)^2 L_{x_0,k}$ would put
$\mathbf{77\%}$ on \texttt{ti=27}. The observed $\mathbf{50\times}$ ratio $L_{\mu,27}/L_{\mu,13}$
is \emph{between} the flat-$L_v$ prediction ($\sim148\times$ from $\dt$ ratio alone, before
$L_v$ shrinkage) and the flat-$L_{x_0}$ prediction ($\sim3100\times$), because both the metric
\emph{and} the actual $L_v$ vary with~$k$.
\item Switching to $L_{x_0}$ (\texttt{xopd\_dk\_space="x0"}) replaces the $\dt_k^2$ weighting with
$\sigma_k^2$ (col.~``\%$L_{x_0}$''), up-weighting high-noise structure steps.
\end{itemize}
This parallels diffusion training, where $\varepsilon$-prediction MSE needs SNR reweighting to avoid
low-noise dominance; here the analogue of ``native'' FM regression is $L_v$, not $L_{xt}$.

\subsection{Cause B --- genuine difficulty near the data manifold (a real signal)}
\label{ssec:causeB}

Near clean, the posterior $\E[z_0\mid x_t]$ sharpens, the velocity field's spatial Jacobian
(Lipschitz constant) grows $\sim 1/\sigma$, and the target encodes sharp, high-frequency, multimodal
detail. Under \texttt{normalize\_d\_k=false} the latent scale is also largest at clean. This part of
the late-step cost is \textbf{real} and tied to final-image quality.

\subsection{The U-shape}
\label{ssec:ushape}

The observation doc's minimum at ti$\approx12$--$13$ and mild high-noise bump at ti$=0$ follow from
$L_{xt}=\dt^2 L_v$ with $L_v$ largest at high noise: early steps have small $\dt^2$ but large $L_v$;
late steps have large $\dt^2$ but shrinking $L_v$; the minimum is where these cross.

\begin{tcolorbox}[colback=red!5,colframe=red!50!black,title=\textbf{Corollary: \texttt{normalize\_d\_k=true} makes it worse, not better}]
\texttt{normalize\_d\_k=true} divides by $2\bar\sigma^2=2\,\mathrm{std\_dev}_t^2\,\lvert\dt\rvert$,
and $\bar\sigma\to0$ toward clean, stacking a \emph{second} amplification on top of $1/\sigma^2$.
This is why transition-KL objectives (DanceGRPO/DDPO-style) intrinsically favor the low-noise
region. \textbf{To equalize timesteps you must go the opposite way.}
\end{tcolorbox}

%% ============================================================
\section{XOPD-specific reading (large teacher)}
\label{sec:xopd}
%% ============================================================

``XOPD benefits more because the teacher is larger'' is correct, but the benefit lives in
\textbf{Cause B}, not A:
\begin{itemize}[leftmargin=1.5em]
\item The 32B$\to$4B teacher's advantage (fine detail, texture, aesthetics) concentrates in the
\emph{low-noise} steps, where on-policy $x_t$ is already near the student's final sample --- the most
valuable place for a big teacher to ``polish''.
\item Structural/semantic correctness (geneval count/position, OCR) is decided at \emph{high-noise}
steps.
\end{itemize}
So ``large teacher'' does \textbf{not} imply ``train only late steps''. It implies ``the late-step
Cause-B component is worth keeping --- but do not mistake the Cause-A amplification for that value.''

%% ============================================================
\section{Unified view: every proposal is a choice of per-step weight $w_k$}
\label{sec:unified}
%% ============================================================

All proposals choose how to weight the three losses (or an equivalent $L_v$) across steps:
$\mathcal{L}=\sum_k w_k L_{\{xt,v,x_0\},k}$.

\begin{center}
\renewcommand{\arraystretch}{1.3}
\footnotesize
\begin{tabular}{@{}p{3.6cm}p{4.5cm}p{5.0cm}@{}}
\toprule
\textbf{direction} & \textbf{loss / weight on $L_v$} & \textbf{effect} \\
\midrule
default & $L_{xt}$, $w_k=1$ $\Rightarrow$ $w_k L_v=\dt_k^2$ & late-step dominance (41\% on ti=27 if $L_v$ flat) \\
1. late-only (DanceOPD) & $L_{xt}$ on $\sigma_k<\tau$ only & polish; drops structure learning \\
2a. $L_{x_0}$ (implemented) & $L_{x_0}$, $w_k=1$ $\Rightarrow$ $w_k L_v=\sigma_k^2$ & up-weights high-noise steps \\
2b. $L_v$ (implemented) & \texttt{xopd\_dk\_space="v"} & native FM velocity MSE \\
2c. EMA balance & $L_{xt} / \mathrm{EMA}(L_{xt})$ & flattens logged scale; may suppress Cause B \\
3. endpoint-sensitivity & weight $\propto\partial z_0/\partial v_k$ & 2a is first-order proxy \\
\bottomrule
\end{tabular}
\end{center}

\subsection{Evaluation}
\label{ssec:eval}

\begin{itemize}[leftmargin=1.5em]
\item \textbf{Direction 1 (late-only)} trains $L_{xt}$ on a late window only; same $\dt_k^2$ factor
within the window. Drops high-noise steps where $L_{x_0}$ (hence structure) matters most.
\item \textbf{Direction 2a ($L_{x_0}$)} replaces $\dt_k^2$ with $\sigma_k^2$ on the shared $L_v$.
\textbf{Direction 2b ($L_v$)} removes the $\dt_k^2$ factor entirely (\S\ref{sec:empirical}).
\textbf{Do not use \texttt{normalize\_d\_k=true}} to equalize (\S\ref{ssec:ushape}).
\item \textbf{Direction 3 (sensitivity)} idealizes weighting $L_v$ by endpoint Jacobian; $L_{x_0}$
is a good first-order proxy.
\end{itemize}

\begin{tcolorbox}[colback=green!5,colframe=green!50!black,title=\textbf{Recommendation}]
Compare \texttt{xopd\_dk\_space} $\in\{\texttt{xt},\texttt{v},\texttt{x0}\}$ on OCR
(\S\ref{sec:experiment}): default \texttt{xt} vs.\ velocity \texttt{v} vs.\ clean-latent
\texttt{x0}, plus late-only \texttt{xt} (direction~1).
\end{tcolorbox}

%% ============================================================
\section{Disambiguating experiment (artifact vs.\ signal)}
\label{sec:experiment}
%% ============================================================

The ablation in \texttt{docs/xopd/progress/2026-07-12-ablation-relay.md} (all on OCR, otherwise
identical) is exactly the test:

\begin{center}
\renewcommand{\arraystretch}{1.3}
\footnotesize
\begin{tabular}{@{}p{3.4cm}p{6.6cm}p{2.6cm}@{}}
\toprule
\textbf{variant} & \textbf{config} & \textbf{reads} \\
\midrule
full-timestep xt-MSE & \texttt{flux2\_klein\_32b\_to\_4b\_l1\_ocr\_1kep.yaml} & default (\texttt{xt}) \\
late-timestep xt-MSE & \texttt{flux2\_klein\_32b\_to\_4b\_l1\_ocr\_selective\_teacher\_1kep.yaml} & direction 1 \\
full-timestep x0-MSE & \texttt{flux2\_klein\_32b\_to\_4b\_l1\_ocr\_xspace\_1kep.yaml} & direction 2a \\
full-timestep v-MSE & \texttt{flux2\_klein\_32b\_to\_4b\_l1\_ocr\_vmse\_1kep.yaml} & direction 2b \\
\bottomrule
\end{tabular}
\end{center}

Interpretation:
\begin{itemize}[leftmargin=1.5em]
\item If \texttt{x0} or \texttt{v} \textbf{improves} high-noise-half learning and structure metrics
\emph{without} degrading clean-end detail $\to$ the original \texttt{xt} dominance was mostly
\textbf{Cause A (artifact)}; reweighting is a net win.
\item If \texttt{x0}/\texttt{v} \textbf{hurts} clean-end detail with insufficient structure gains
$\to$ the late-step \textbf{Cause B (real value)} dominates; consider soft late-window \texttt{xt}.
\end{itemize}

\begin{remark*}
Logged \texttt{train/d\_k/\{ti\}} magnitudes differ across \texttt{v}/\texttt{xt}/\texttt{x0}; compare
per-timestep \emph{shape} and downstream eval (\S\ref{sec:empirical}), not raw values across spaces.
\end{remark*}

\end{document}
